\chapter{Reconstruction of quantum dynamics through coherent amplitudes over route genealogies}

\epigraphtext{The wave is not added to the history: it is the linear form in which the complete history simultaneously preserves all its possible paths.}{HMT--MD propagation principle}

Measurement has already been defined as a coupling among the system, the apparatus, and the reader. It is now necessary to explain what propagates before the result is published, how one and the same structure supports both a wave and a persistent trajectory, and how observable locality coexists with a global genealogy.

The response of Triadic Modular Holography begins from a distinction that the pointwise formulation tends to obscure. A state is not exhausted by the cell in which it lies or by the real number ultimately read from it. It also preserves the route, orientation, quotients, carries, boundary, and return that brought it there. Upon linearizing this space of histories, the sum over paths does not appear as a metaphor added a posteriori: it is the exact expansion of the propagator. The wave is the coherent evaluation of the totality of routes; in the following chapter, the particle will be the section that persists along one of them.

\section{The global amplitude of \(1\) family of paths}

Let \(X\) be a finite set of observable states at a given depth, and let
\(U:\mathbb C^X\to\mathbb C^X\) be the one-step propagation operator. A route of
length \(R\) from \(i\) to \(f\) is a sequence

\[
  \gamma=(x_0=i,x_1,\ldots,x_R=f)
\]

of admissible transitions. Its ordered weight is

\[
  w_U(\gamma)=\prod_{t=1}^{R}U_{x_t x_{t-1}}.
\]

Although scalar weights commute, the indexed sequence of transitions matters:
one is not counting only how many routes reach \(f\), but which factor each
step assigns to that history. Multiplication composes the transitions along a route;
addition gathers alternative histories that terminate in the same state. For
operator-valued weights, the temporal product is moreover genuinely noncommutative. This
difference between a product within a route and a sum among routes is the linear
version of the double reading that runs throughout HMT.

\begin{theorem}[Finite path sum]
For all \(R\geq 1\),
\[
  (U^R)_{fi}
  =
  \sum_{\substack{\gamma:i\to f\\|\gamma|=R}}
  w_U(\gamma).
\]
If \(\psi_0\in\mathbb C^X\) is the initial amplitude, then
\[
  \psi_R(f)
  =
  \sum_{i\in X}
  \sum_{\substack{\gamma:i\to f\\|\gamma|=R}}
  w_U(\gamma)\,\psi_0(i).
\]
\end{theorem}

\begin{proof}
The assertion for \(R=1\) is the definition of \(U\). When an expansion
of \(U^R\) is multiplied by \(U\), every route of length \(R\) receives exactly one final
admissible edge. Distributivity sums all possible extensions, and
multiplication preserves the order of their weights. By induction, every route of length
\(R+1\) appears exactly once. The second equality applies the first to
\(\psi_0\).
\end{proof}

This identity has considerable conceptual content. A component
\(\psi_R(f)\) is not an isolated property of the point \(f\). It is the contraction of the
complete genealogy of routes reaching it. Two histories may end at the
same visible value and nevertheless contribute opposite phases; they may also
reinforce one another. Interference consists precisely in summing amplitudes before
forming a positive observable. There is no need to introduce an additional material wave
traversing a space distinct from that of the histories: the wave structure
resides in the linear composition of the routes.

When \(U\) is unitary,

\[
  \|\psi_R\|_2=\|U^R\psi_0\|_2=\|\psi_0\|_2.
\]

The total quantity that may subsequently be normalized as a probability remains
constant. If the physical reader assigns to an output \(f\) the weight

\[
  p_R(f)=\frac{|\psi_R(f)|^2}{\|\psi_R\|_2^2},
\]

then \(p_R\) is a normalized distribution, and its evolution is compatible with unitary conservation. This is the precise point at which route amplitude becomes outcome statistics. The chapter on measure has already located the coupling that publishes an output; here, the object evaluated by that coupling has been identified.

The construction also clarifies why a visible number is not sufficient to reconstruct
the physics. The map

\[
  \{\text{deep routes}\}\longrightarrow X
\]

generally forgets the relative phase between paths. Knowing only the
final distribution of cells does not permit inversion of the sum or recovery of the
genealogy. The complete HMT state preserves what the screen contracts. Hence
memory is not an interpretive ornament: it is the datum on which interference depends.

At the refinement limit, the correct question is not whether the discrete route ``becomes'' mysteriously continuous, but whether the propagators form a compatible family. For refinement projections \(p_{m+1,m}:X_{m+1}\to X_m\), coherence requires the cylindrical observables and the transported states to agree. When this compatibility is exact, the amplitude of a coarse history is the sum of the amplitudes of all its fine prolongations. When the defects are summable, the expectations of uniformly bounded cylindrical observables form a Cauchy family. Operational continuity is thereby obtained without erasing the discrete structure that generates it.

\begin{principle}[The wave as linear memory]
The finite-depth HMT wave function is the linear evaluation of the route
extension. It coherently preserves the alternatives compatible with the
preparation and the boundary; the observed value is a subsequent projection
of that global amplitude.
\end{principle}

This formulation is stronger than asserting that ``everything occurs along many paths''.
It specifies the domain, the weight of each path, the composition rule, and the
preserved norm. It also contains its negative control: if a purported amplitude does not
respect route composition, counts a history twice, consults the
future result in order to assign the weight, or loses the norm under a propagator
declared unitary, it is not the wave function constructed here.

\section{Pilot wave, route performed and double orientation}

The preceding structure already contains the two components that a pilot-wave theory
must distinguish. The amplitude \(\psi_t\) globally organizes the alternatives;
a persistent route \(x_t\) records the history actually realized. The wave
and the particle are not two rival substances. They are two levels of reading of the same
enriched state: the linear field of possibilities and the section that preserves its
identity while traversing it.

The guiding law must convert this kinematic architecture into dynamics. Its
minimal form is a map

\[
  G:(\psi_t,x_t)\longmapsto x_{t+1},
\]

o, if the update is stochastic, a kernel
\[
  P_{\psi_t}(y\mid x).
\]

The central condition is not verbal, but mathematical: the guided distribution must be
equivariant with amplitude evolution,

\[
  |\psi_{t+1}(y)|^2
  =
  \sum_x P_{\psi_t}(y\mid x)\,|\psi_t(x)|^2.
\]

This equality ensures that a population initially distributed according to
\(|\psi_t|^2\) remains so distributed after one step. The guide cannot be
chosen retrospectively to reproduce each detector; it must be fixed by
the state, propagator, orientation, and register before the output is known.
HMT adds a possibility that a pointwise description does not possess: \(G\) may
depend on route and memory cocycles that do not descend to the visible value
\(x_t\). The local action can thereby be informed by a global genealogy without
the appearance of a supplementary force traveling from the future.

The bidirectional reading further develops this structure. For an oriented route
\(\gamma\) with

\[
  \partial\gamma=x_{\rm abs}-x_{\rm em},
\]

the return involution defines

\[
  \gamma^\vee=C_{\rm rev}\gamma,
  \qquad
  \partial\gamma^\vee=x_{\rm em}-x_{\rm abs}.
\]

These are not two independently and freely chosen paths. They are the two
orientations of the same history. The internal length is invariant,

\[
  L_{\rm int}(\gamma^\vee)=L_{\rm int}(\gamma),
  \qquad
  L_{\leftrightarrow}=2L_{\rm int}(\gamma),
\]

and inversion interchanges the conjugate channels. Preparation and absorption
cease to be ontologically disconnected endpoints: they belong to a closure in
which the outward and return paths must be compatible.

When this involution is represented in the physical theory of propagators, the two orientations admit retarded and advanced readings. The retarded component transports the preparation toward absorption; the advanced component transports the return condition toward emission. The observable history is their mutual adjustment. This formulation precisely expresses the intuition of a pilot wave with advance and retardation: the local guide receives information from both boundaries of the same oriented object.

This provides no authorization to send controllable messages into the past.
An advanced solution of an equation and a signalable future choice are distinct
objects. The physical representation must prove that the propagators satisfy
the evolution equation, that the boundary conditions compose
consistently, and that the marginal probabilities do not permit modulation of an
earlier output by a later decision. Dual orientation is the
support; observable causality is a property that must be verified on its
physical reader.

The correlation between the two sheets can be formalized in two Hilbert spaces,
\(\mathcal H_{\rm adv}\) and \(\mathcal H_{\rm ret}\). For a distribution
\(p=(p_i)\), the state

\[
  |\Psi_p\rangle
  =
  \sum_i\sqrt{p_i}\,
  |i\rangle_{\rm adv}\otimes|i\rangle_{\rm ret}
\]

has reduced matrices with spectrum \((p_i)\) and entropy

\[
  S_{\rm ent}(p)=-\sum_i p_i\log p_i.
\]

The equality is exact once the bilayer factorization has been declared. A single history
has Schmidt rank one and zero entropy; several correlated histories
produce positive entropy. The result does not simply identify this entropy with
\(\log\phiHMT\), with the information in a triadic cell, or with thermodynamic
entropy. Each belongs to a different object. What it proves is that
outbound and return can form a nonseparable state and that the shared information
is calculated without metaphors.

\begin{test}[Bidirectional guidance]
A physical realization of the HMT pilot wave must fix, before comparison, the propagator \(U\), the law \(G\) or the kernel \(P_\psi\), the involution \(C_{\rm rev}\), and the boundary conditions. It is falsified if it loses equivariance, assigns two incompatible routes to the same complete state, needs to know the future detector in order to choose the trajectory, or permits controllable signaling against the published causal orientation.
\end{test}

This proof does not relegate the proposal to a mere possibility. On the contrary, it shows that the mathematical core has already been constructed and identifies the sole locus at which physical dynamics must act. The sum over paths is not pending; the persistent route is not pending; nor is inversion. What is decided experimentally is which Hamiltonian realizes the guide and which boundary conditions turn the two sheets into nature's advanced and retarded waves.

\section{Bell: locality of the data and globality of the genealogy}

The distinction between complete state and projected value permits the Bell
problem to be formulated precisely. Consider settings \(a,b\in\{0,1\}\), outcomes
\(A,B\in\{-1,+1\}\), and a shared state \(\lambda\). A factorizable
local theory satisfies

\[
  P(A,B\mid a,b)
  =
  \int_\Lambda
  P_A(A\mid a,\lambda)\,
  P_B(B\mid b,\lambda)\,
  \dd\rho(\lambda),
\]

with a measure independent of the adjustments,
\[
  \dd\rho(\lambda\mid a,b)=\dd\rho(\lambda).
\]

Defining
\[
  E_{ab}
  =
  \sum_{A,B=\pm1}AB\,P(A,B\mid a,b),
\]

we obtain
\[
  S_{\rm CHSH}=E_{00}+E_{01}+E_{10}-E_{11},
  \qquad |S_{\rm CHSH}|\leq2.
\]

HMT preserves this theorem. Adding hidden memory to \(\lambda\) does not violate the bound if
factorization, setting independence, and the absence of
postselection are maintained. This control is important: it prevents the richness of TPK from being turned into
an automatic explanation of every correlation.

The holographic contribution is formulated differently. The two wings of an experiment may be local projections of the same enriched section, with a common boundary and genealogy. The datum published on each wing is local; the state from which both descend need not factor into two independent deep histories. Nonlocality then consists not in a signal crossing space after the settings have been chosen, but in the impossibility of decomposing the global preparation into two autonomous complete states.

Non-signaling and Bell locality are not equivalent. It may be satisfied

\[
  \sum_B P(A,B\mid a,b)=P(A\mid a),
  \qquad
  \sum_A P(A,B\mid a,b)=P(B\mid b),
\]

and nevertheless fail the preceding factorization. This is the precise domain that an HMT realization must occupy: global correlation without a superluminal communication channel. The contextual reader may depend on the setting in how it projects the common state, but an observer's marginal distribution cannot depend on the remote choice.

\begin{construction}[HMT Protocol for Bell]
An enriched state \(x\), a preparation measure, and four
readers \(L_{ab}\) are fixed first. A binary outcome map is then declared, and
the complete table is calculated
\[
  P_{\rm HMT}(A,B\mid a,b).
\]
The construction is physically valid only if it is frozen before the data are
consulted, is normalized, has no hidden postselection, and publishes simultaneously
\(S_{\rm CHSH}\) and the four marginal no-signaling controls.
\end{construction}

The so-called Algorithm 3 belongs to this specific construction. Its purpose is not to prove the classical theorem again, but to construct the descent from the shared genealogy to a contextual distribution. If it preserves all the Bell premises, it will necessarily produce \(|S|\leq2\). If it produces \(|S|>2\), it must show explicitly which factorization fails to descend and why no-signaling is nevertheless preserved. This is the complete holographic explanation of the locality problem: the observable space may be separable even though the history generating it is not separable.

The proposal is especially natural in a theory in which sum and product are
distinct readers. Bell factorization seeks to express a joint probability
as a conditioned product of local responses; quantum amplitude, by contrast,
arises by summing histories before projecting them. If projection is
performed too early, the common phase is lost and only the local polytope remains.
If the global state is preserved until the joint reading, interference terms
survive that do not admit such a decomposition. The incompatibility is not a
license to violate logic: it is the typed difference between composing alternatives
by sum and composing local probabilities by product.

\section{From the path theorem to quantum contrast}

The resulting quantum architecture is articulated through a single chain of maps. HMT
constructs an enriched discrete space of histories; its linearization produces the
coherent sum over paths; unitarity preserves the norm; a persistent section
makes it possible to speak of a trajectory; oriented inversion yields the outward-and-
return pair; two-sheet factorization quantifies its correlation; and Bell analysis
distinguishes locality of the datum from separability of the genealogy. These are not
juxtaposed analogies, but successive maps on the same state.

The theory thus offers a unified reading of quantum features.
Superposition is the linear coexistence of prolongations; interference is their
phased sum; the particle is the persistence of a section; apparent reduction
is the selection of a fiber by the reader; entanglement is the
nonfactorization of sheets or subsystems; and nonlocality is the globality of the
genealogy preserved under locally nonsignaling projections.

Physical tests are added not to weaken this reading, but to determine its natural realization. A finite interferometric network can implement \(U\) and compare the path expansion with the final pattern. A reconstruction of currents can test the guidance law. Controlled inversion of the boundary conditions can verify the reciprocity \(\gamma\leftrightarrow\gamma^\vee\). A Bell experiment with sealed readers can determine whether deep memory generates a nonfactorizable, no-signaling distribution.

\begin{test}[Joint Quantum Compatibility Certificate]
The certificate distinguishes four independent validity conditions:
\begin{enumerate}
  \item the propagator not reproduce by coherent sum the amplitudes observadas;
  \item no natural guiding law is equivariant with \(|\psi|^2\);
  \item the double orientation does not admit reciprocal physical propagators without
  improper causal signaling;
  \item the Bell descent is locally factorizable or violates no-signaling.
\end{enumerate}
Failure of an interface localizes the incorrect map; it does not erase the finite
route theorems that precede it.
\end{test}

The positive thesis is therefore clear. Quantum theory does not require the
history of a system to be abandoned; it requires what is meant by history to be enriched. A point
without genealogy cannot by itself carry interference, return, or correlation.
An oriented extension can do so. Its wave is the coherent memory of all
compatible paths; its particle is the closure that persists. The passage between
the two forms is not an ontological leap: it is the passage from the linear totality of routes to
a realized section of that same totality.
